\emph{Cet exercice est un questionnaire à choix multiple.
Pour chaque question, une seule des quatre réponses proposées est exacte. Le candidat indiquera sur sa copie le numéro de la question et la réponse choisie. Aucune justification n'est demandée. Une réponse fausse, une réponse multiple ou l'absence de réponse à une question ne rapporte ni n'enlève de point.}

Les quatre questions sont indépendantes.
 

Pour chaque question, une seule réponse correcte. Aucune justification n’est demandée.

Dans tout l'exercice, on considère que l'espace est muni d'un repère orthonormé

\Oijk.

On considère:

\begin{itemize}[label=$\bullet~$]
\item les points A$(-3~;~1~;~4)$ et B(1~;~5~;~2)
\item le plan $\mathcal{P}$ d'équation cartésienne $4x + 4y -2z + 3 = 0$
\item la droite $(d)$ dont une représentation paramétrique est $\left\{\begin{array}{l c l}
x&=&- 6 + 3t\\y&=&\phantom{-}1\\z&=&\phantom{-}9 - 5t
\end{array}\right.,$ où $t \in~\R$.
\end{itemize}

\medskip
\begin{enumerate}
\item Les droites (AB) et $(d)$ sont:
	\begin{enumerate}
		\item sécantes non perpendiculaires. 
		\item  perpendiculaires.
		\item non coplanaires.
		\item parallèles.
	\end{enumerate}
\item La droite (AB) est:
	\begin{enumerate}
		\item incluse dans le plan $\mathcal{P}$.
		\item strictement parallèle au plan $\mathcal{P}$.
		\item sécante et non orthogonale au plan $\mathcal{P}$.
		\item orthogonale au plan $\mathcal{P}$.
	\end{enumerate}
\item On considère le plan P' d'équation cartésienne $2x+y+6z+5 = 0$. 

Les plans $\mathcal{P}$ et $\mathcal{P}'$ sont:
	\begin{enumerate}
		\item sécants et non perpendiculaires.
		\item perpendiculaires.
		\item confondus.
		\item strictement parallèles.
	\end{enumerate}
\item On considère le point C$(0~;~1~;~-1)$. La valeur de l'angle $\widehat{\text{BAC}}$ arrondie au degré est: 
	\begin{enumerate}
		\item 90\ensuremath{^\circ}
		\item 51\ensuremath{^\circ}
		\item 39\ensuremath{^\circ}
		\item 0\ensuremath{^\circ}
	\end{enumerate}
\end{enumerate}

\bigskip
